\ansclass{4}{Traffic \& Edge}
\ans{1} (a) \final{Cold $\approx$ 130 ms, warm $\approx$ 80 ms}: connection reuse removes about a third of the cold cost; both are under a 200 ms p99 target. (b) At 10 ms: \final{40, 30, 20 and 10 ms}; at 150 ms: \final{600, 450, 300 and 150 ms}. Distance dominates: with a same-city edge even TLS 1.2 costs only 40 ms. (c) 0-RTT data \final{can be replayed} by an attacker; only idempotent requests may use it.
\careful{optimizing the server's 40 ms and ignoring the handshakes, which cost more as soon as the user is far away.}
\ans{2} (a) One query per resolver per TTL: $5\,000 / 60 \approx$ \final{80 queries/s}, nothing like 2 M. (b) $30 \times 3 + 60 =$ \final{150 s $\approx$ 2.5 min}, best case: some resolvers ignore low TTLs and open connections stay on the dead region. (c) \final{Lower the TTL to 60 s on Monday} (at least one old TTL before the move), so Tuesday's change spreads in a minute instead of an hour. (d) A load balancer checks every few seconds and stops routing at once; DNS answers sit in caches for the TTL.
\careful{assuming DNS failover is instant. It is detection $+$ TTL, minutes in the best case.}
\ans{3} (a) \final{20\,000, 2\,000 and 200 req/s}: each nine divides the origin load by 10. (b) $10\,000 \to 2\,000$ req/s: \final{5$\times$ fewer}, about 80\% of the origin servers saved. (c) $2 \times 10^6 / 5 =$ \final{400\,000 req/s} at the edge; with \texttt{s-maxage=1}, each PoP asks at most once a second: \final{$\approx$ 50 req/s at the origin}, data \final{at most 1 s old}; with a shield, $\approx$ \final{1 req/s}. (d) Bundle: \texttt{public, max-age=31536000, immutable}; event HTML: \texttt{public, max-age=60, s-maxage=300, stale-while-revalidate=60}; seat JSON: \texttt{public, max-age=0, s-maxage=1}; order and payment: \final{\texttt{private, no-store}}.
\careful{thinking in hit ratio instead of miss ratio: from 95\% to 99\% the origin does not get 4\% lighter, it gets 5$\times$ lighter.}
\ans{4} (a) $5\,000/800 = 6.25 \to$ \final{7 at 100\%}; run \final{8 or more} (headroom and a spare, class 5). (b) 30 $\to$ \final{B}, 50 $\to$ \final{C}, 90 $\to$ \final{A} (wraps around). After D: only 50 moves, \final{to D}; in general only \final{positions 41--55} (between B and D) move. (c) $\bmod N$: about 5/6, \final{$\approx$ 33\,000 keys}; ring: about 1/6, \final{$\approx$ 6\,700 keys}. (d) Random points split the ring into unequal arcs; \final{100--200 virtual nodes} per server average the arcs out. (e) Up to $5 \times 2 =$ \final{10 s} of errors for its share of traffic. Draining: on shutdown, readiness fails first, the pod stops receiving new requests and finishes the in-flight ones (up to $\approx$ 30 s), so a deploy causes no 502s.
\careful{believing consistent hashing fixes a hot key: the stadium seat map is one key and still lands on one node.}
\ans{5} (a) \final{20 accepted, 30 get 429} with \texttt{Retry-After}. (b) The bucket loses $14 - 10 = 4$ tokens/s: everything is accepted for $20/4 =$ \final{$\approx$ 5 s}, then \final{exactly 10 req/s} (4 req/s rejected). (c) $20/10 =$ \final{2 s}. (d) \final{50 accepted, 50 rejected}; refill $50/5 =$ \final{10 s}; long-run average \final{$\le$ 5 req/s}. Capacity is the burst; the rate is the sustained limit.
\careful{reading the capacity as the per-second limit: a bucket of 20 refilled at 10/s accepts 20 in the first instant but only 10/s after that.}
\ans{6} (a) \final{200 in about 2 s}: 100 at 10:00:59 and 100 more at 10:01:00, when the counter resets. (b) $30 + 80 \times (1 - 15/60) = 30 + 60 =$ \final{90 $<$ 100: allowed}. (c) Up to $6 \times 10 =$ \final{60 req/s}: each replica sees a sixth of the traffic. Fix: \final{one shared atomic counter} (Redis $+$ Lua script), or limit/replicas locally as an approximation. (d) $50\,000 \times 100 = 5$ M/min $\approx$ \final{83\,000 req/s}, all accepted. Better keys: \final{account, card, device fingerprint}, plus bot detection and the waiting room.
\careful{limiting only by IP: botnets have thousands of IPs, and thousands of real fans share one IP behind carrier NAT.}
\ans{7} (a) $2 \times 10^6 \times 0.5 =$ \final{1 M req/s}, almost all ``nothing changed''. Push with \final{SSE} (server-sent events) during the sale: one event per change replaces millions of empty answers; the cost is 2 M open connections (tens of servers, a pub/sub behind them, a plan for reconnection storms). (b) A good answer: DNS TTL 60 s with latency routing and failover; CDN with versioned assets for a year, event page \texttt{s-maxage=300}, seat JSON micro-cached 1 s, holds and orders never cached; WAF with managed rules, bot detection and $\le$ 5 \texttt{POST /login} per IP per minute; gateway \final{10 req/s, burst 20 per fan} on \texttt{/api/**} and \final{1 req/s, burst 3} on \texttt{POST /api/holds}; a waiting room admitting \final{$\approx$ 2\,000 fans per minute}; \texttt{inventory} allows 6 seats per hold. Each hop removes traffic: the core sees thousands of req/s instead of millions. (c) Decide in advance: BoxOffice \final{fails open with a local limit of 1/6 of the global one} (6 replicas), so a Redis outage neither blocks fans nor removes all protection.
\careful{drawing the edge as one box: each hop must say what it removes and what happens when it fails.}
